%% -*- mode: TeX -*-
%% -*- coding: utf-8 -*-
\documentclass[12pt,reqno]{amsart}

% WinEdt/pdflatex UTF-8 support (safe for pdfLaTeX)
\usepackage[utf8]{inputenc}

% =====================================================================
% Template for writing math papers (amsart + mathpaper.sty).
%
% Usage:
% - Edit the "Paper metadata" block below.
% - Put bib entries into references.bib (or change \bibliography{...}).
% - Compile: latexmk -pdf main.tex
% =====================================================================

\usepackage{mathpaper}

% =====================================================================
% Paper metadata (EDIT THIS BLOCK)
% =====================================================================
\title[Short Title]{Full Paper Title Goes Here}
\date{\today}
\keywords{keyword 1; keyword 2; keyword 3}
\subjclass[2020]{Primary 00A05; Secondary 00A06}
%\thanks{Optional funding / acknowledgements footnote.}

\author{First Author}
\address{Department, University, City, Country}
\email{author1@example.edu}

%\author{Second Author}
%\address{Department, University, City, Country}
%\email{author2@example.edu}

\begin{document}
\begin{abstract}
Write a concise abstract (typically 100--200 words) summarizing the main results
and methods. Avoid undefined notation when possible.
\end{abstract}
\maketitle

\section{Introduction}\label{sec:intro}
Explain the problem, motivation, and main contributions. A typical outline:
\begin{itemize}
  \item Background and related work, e.g. \cite{Doe2024Example}.
  \item arXiv preprints can be cited too, e.g. \cite{Smith2025ArxivExample}.
  \item Statement of main results.
  \item Organization of the paper.
\end{itemize}

\subsection{Main results}
\begin{thm}[Example main theorem]\label{thm:main}
Let \(X\) be a set. Then \(X=X\).
\end{thm}

\begin{rmk}
Use \verb|\Cref{thm:main}| for cross references (\Cref{thm:main}).
\end{rmk}

\section{Preliminaries}\label{sec:prelim}
Define notation and recall standard facts.

\begin{defn}[Example definition]
A \emph{widget} is a thing satisfying property~(W).
\end{defn}

\begin{lem}[Example lemma]\label{lem:tool}
If \(a=b\), then \(a^2=b^2\).
\end{lem}

\begin{proof}
\step[Idea]
Square both sides.
\step[Finish]
The conclusion follows.
\end{proof}

\section{Proof of the main theorem}\label{sec:proof}
\begin{proof}[Proof of \Cref{thm:main}]
Combine \Cref{lem:tool} with the definitions.
\end{proof}

\appendix
\section{An optional appendix}\label{app:extra}
Put technical lemmas or computations here.

\section*{Acknowledgements}
Write acknowledgements here (optional).

\bigskip
\bibliography{references}
\bigskip
\end{document}

